\renewcommand{\thesection}{S\arabic{section}}
\setcounter{section}{0}

\section{Construction of the illustrative example in Figure 1}
\label{sec:Appendix_S1}

Figure~\ref{fig:intro} provides a data-informed illustration of how
delay filtering suppresses rapid differences between candidate upstream
epidemic trajectories. The figure is intended to illustrate the
observation and separability framework rather than to provide a
substantive reconstruction of the U.S. epidemic.

\subsection{Observed downstream incidence}

Daily U.S. confirmed-case incidence was obtained from the Johns Hopkins University Center for Systems Science and Engineering COVID-19 data stream \cite{dong2020interactive}, accessed through the Delphi Epidata API \cite{farrow2015delphi}. We used the national-level \texttt{jhu-csse/confirmed\_incidence\_num} signal from April 1, 2020, through June 30, 2021. The series was aligned to a complete daily sequence, with missing values linearly interpolated. Figure~\ref{fig:intro}

displays the period from September 1, 2020, through March 31, 2021.

\subsection{Smoothed downstream mean and observation noise}

To separate broad epidemic variation from day-to-day observational
fluctuation, we estimated a smooth downstream mean
\(\widehat{\mu}(t)\) using second-order trend filtering. The effective
degrees of freedom were selected over a prespecified grid. Conditional
on each candidate smooth mean, Gaussian, Poisson, and negative-binomial
observation models were fitted, and the combination of trend-filter
complexity and observation model was selected by the Akaike information
criterion. The resulting residual or model-based variance scale was used
to define the effective observation-noise level shown in
Fig.~\ref{fig:intro}F.

\subsection{Delay distribution}

The delay kernel was based on the published COVID-19
symptom-onset-to-case-report distribution reported by
Zhang et al.~\cite{zhang2020evolving}. A lognormal representation was
converted to a discrete daily probability mass function by integrating
the fitted distribution over consecutive one-day intervals. The kernel
was truncated at 60 days and renormalized to sum to one.

\subsection{Construction of the reference upstream trajectory}

The reference upstream trajectory \(U_0(t)\) was obtained by solving a
nonnegative regularized inverse problem. Let \(K_g\) denote the linear
convolution operator induced by the delay kernel \(g\). We estimated
\(U_1\) by minimizing
\[
\left\|K_g U-\widehat{\mu}\right\|_2^2
+
\lambda\left\|\Delta^2 U\right\|_1
\]
subject to \(U(t)\geq 0\), where \(\Delta^2\) is the second-difference
operator. The regularization term stabilizes the inverse problem and
produces a smooth reference trajectory whose delayed projection follows
the estimated downstream mean.

\subsection{Construction of the alternative upstream trajectory}

The alternative trajectory \(U_1(t)\) was constructed to differ visibly
from \(U_0(t)\) at short temporal scales while producing a similar
downstream projection. Candidate perturbations were formed from mixtures
of sinusoidal components concentrated near frequencies corresponding to
7- and 3-day periods. For each candidate perturbation \(h(t)\), we
considered
\[
U_1(t)=\max\{U_0(t)+a h(t),0\},
\]
where the amplitude \(a\) was increased to the largest value satisfying
both
\[
\operatorname{RMSE}(K_gU_1,K_gU_0)
\leq
\gamma\,\operatorname{RMSE}(K_gU_0,y)
\]
and
\[
\operatorname{RMSE}(K_gU_1,y)
\leq
\eta\,\operatorname{RMSE}(K_gU_0,y).
\]
For the displayed example, \(\gamma=1\) and \(\eta=1.05\). Among feasible
candidates, we selected the perturbation producing the largest visible
upstream separation while retaining substantial spectral energy near
the two target frequencies.

\subsection{Spectral quantities}

The upstream spectral contrast shown in Fig.~\ref{fig:intro}D is
\[
\left|\widehat{U}_0(f)-\widehat{U}_1(f)\right|^2.
\]
The delay attenuation in Fig.~\ref{fig:intro}E is
\[
\left|\widehat{G}(f)\right|^2,
\]
and the filtered contrast in Fig.~\ref{fig:intro}F is
\[
p^2
\left|\widehat{G}(f)\right|^2
\left|\widehat{U}_0(f)-\widehat{U}_1(f)\right|^2.
\]
The horizontal reference line denotes the effective observation-noise
level under the selected observation model. These quantities are shown
to illustrate the frequency-wise components that motivate the
separability functional introduced in the main text.